% English twin of the French section
\section{Coordinate geometry: vectors and lines, a synthesis}

\fichehead{Tue~30~Mar.~2027}{Vectors and collinearity: consolidation}
\begin{fiche}
\qitem{Warm-up · Quick}{GEO}{$\vec u\vc{3}{-2}$ and $\vec v\vc{-1}{5}$. Calculate $\vec u+\vec v$ and $2\vec u$.}{$\vc{2}{3}$; $\vc{6}{-4}$.}{}
\qitem{Core · Method}{GEO}{Are the points $A(-1;2)$, $B(1;3)$ and $C(5;5)$ collinear?}{$\vect{AB}\vc21$, $\vect{AC}\vc63$: $2\times3-1\times6=0$: collinear.}{}
\qitem{Core + · Method}{GEO}{$A(-1;2)$, $B(3;1)$, $C(4;4)$. Find $D$ such that $ABCD$ is a parallelogram.}{$\vect{AB}=\vect{DC}$: $D=A+\vect{BC}=(-1+1;2+3)=(0;5)$.}{}
\qitem{Review · Quick}{PRO}{In a class of $24$ pupils, $18$ have school lunch. What proportion is this, as a percentage?}{$\dfrac{18}{24}=0.75=75\,\%$.}{}
\qitem{Thread · Prove}{LIT}{Show that $(n+1)^2-(n-1)^2=4n$ for every integer $n$. What can be concluded about the difference of the squares of two consecutive odd numbers ($n$ even)?}{$(n^2+2n+1)-(n^2-2n+1)=4n$. For $n=2k$ even, the difference of the squares of two consecutive odd numbers is $8k$: a multiple of $8$.}{The two integers $n-1$ and $n+1$ are odd when $n$ is even.}
\end{fiche}

\fichehead{Thu~1\textsuperscript{st}~Apr.~2027}{Determinant and equations of lines: consolidation}
\begin{fiche}
\qitem{Warm-up · Quick}{GEO}{What is the gradient of the line $3x+2y-6=0$?}{$y=-\dfrac32x+3$: gradient $-\dfrac32$.}{}
\qitem{Core · Method}{GEO}{Equation of the line $(AB)$ with $A(1;2)$ and $B(4;8)$.}{Gradient $\dfrac{8-2}{4-1}=2$; $y=2x+p$, $2=2+p$: $y=2x$; or $2x-y=0$.}{}
\qitem{Core + · Method}{GEO}{Are the lines $2x-y+1=0$ and $4x-2y+5=0$ parallel, coincident or intersecting?}{Direction vectors $\vc12$ and $\vc24$ are collinear: parallel; $(0;1)$ satisfies the first but $4\times0-2+5=3\neq0$: strictly parallel.}{}
\qitem{Review · Quick}{PUI}{Calculate $10^{-2}\times5\times10^4$.}{$5\times10^{2}=500$.}{}
\qitem{Thread · Method}{EQ}{Solve in $\R$: $7-2(x-1)=3x-6$.}{$x=3$}{Pitfall: $-2(x-1)=-2x-2$ (sign error).}
\end{fiche}

\fichehead{Mon~5~Apr.~2027}{Plane geometry problems: choosing a method}
\begin{fiche}
\qitem{Warm-up · Quick}{GEO}{Distance $AB$ for $A(0;0)$ and $B(5;12)$?}{$\sqrt{25+144}=13$.}{}
\qitem{Core · Method}{DEM}{$A(0;0)$, $B(4;2)$, $C(-2;4)$. What type of triangle is $ABC$?}{$AB^2=20$, $AC^2=20$, $BC^2=36+4=40=AB^2+AC^2$: right-angled isosceles at $A$.}{}
\qitem{Core + · Method}{DEM}{Give two methods to prove, using coordinates, that $ABCD$ is a parallelogram.}{(i) $\vect{AB}=\vect{DC}$; (ii) the diagonals $[AC]$ and $[BD]$ have the same midpoint.}{}
\qitem{Review · Reading}{LEC}{In a class of $30$ pupils, a bar chart shows $x$ pupils with the mark $12$ and $2x$ pupils with the mark $14$; the other marks account for $18$ pupils. Find $x$.}{$x+2x+18=30$, so $3x=12$ and $x=4$.}{}
\qitem{Thread · Quick}{SOL}{A pyramid has a square base of side $6$ cm and height $5$ cm. Calculate its volume.}{$V=\dfrac13\times6^2\times5=60$ cm$^3$.}{}
\end{fiche}

\fichehead{Tue~6~Apr.~2027}{Situations from other subjects}
\begin{fiche}
\qitem{Warm-up · Quick}{MOD}{A moving object has position $x(t)=3t+2$ (in m, $t$ in s). Where is it at $t=4$?}{$x(4)=14$ m.}{}
\qitem{Core · Method}{MOD}{A second object follows $x_2(t)=5t-4$. At what time and where do the objects meet?}{$3t+2=5t-4\Leftrightarrow t=3$ s; position $x=11$ m.}{}
\qitem{Core + · Quick}{GEO}{A velocity vector has coordinates $\vc34$ (in m/s). What is its magnitude?}{$\sqrt{3^2+4^2}=5$ m/s.}{}
\qitem{Review · Quick}{CAL}{Calculate $\dfrac12-\dfrac35$.}{$\dfrac{5}{10}-\dfrac{6}{10}=-\dfrac1{10}$.}{}
\qitem{Thread · Write-up}{DEM}{$ABC$ is isosceles at $A$ and $H$ is the foot of the altitude from $A$. What can be stated about $H$?}{$H$ is the midpoint of $[BC]$ (the altitude from the apex is also the perpendicular bisector of the base).}{}
\end{fiche}

\fichehead{Thu~8~Apr.~2027}{Communicating a proof: implication, converse, equivalence}
\begin{fiche}
\qitem{Warm-up · Quick}{DEM}{State the converse of: ``if a quadrilateral is a rhombus, then its diagonals are perpendicular''.}{``If a quadrilateral has perpendicular diagonals, then it is a rhombus''.}{}
\qitem{Core · True/false}{DEM}{True or false? Justify. The previous converse is true.}{False. Counter-example: a kite (not a rhombus) has perpendicular diagonals.}{}
\qitem{Core + · Write-up}{DEM}{Is the equivalence ``$x^2=4\Leftrightarrow x=2$'' correct? Correct it.}{No: $x=-2$ also works. We have the implication $x=2\Rightarrow x^2=4$ and the equivalence $x^2=4\Leftrightarrow(x=2\text{ or }x=-2)$.}{}
\qitem{Review · Quick}{CMP}{Compare $2^{10}$ and $10^3$.}{$2^{10}=1\,024>1\,000=10^3$.}{}
\qitem{Thread · Quick}{REP}{Between which consecutive integers does $\dfrac{7}{4}$ lie on a number line? Give its decimal value.}{Between $1$ and $2$; $1.75$.}{}
\end{fiche}

\fichehead{Mon~26~Apr.~2027}{Class discussion: comparing methods}
\begin{fiche}
\qitem{Warm-up · Quick}{GEO}{Distance between $A(-2;1)$ and $B(1;5)$?}{$\sqrt{3^2+4^2}=5$.}{}
\qitem{Core · Method}{GEO}{Give three methods to prove that $A$, $B$, $C$ are collinear.}{(i) $\deter{\vect{AB}}{\vect{AC}}=0$; (ii) the equation of $(AB)$ is satisfied by $C$; (iii) $AB+BC=AC$: equality of distances, equivalent to $B\in[AC]$.}{}
\qitem{Core + · Reflection}{GEO}{To prove that two lines given by $y=m_1x+p_1$ and $y=m_2x+p_2$ are parallel, which method is quickest?}{Compare the gradients: $m_1=m_2$; strictly parallel if in addition $p_1\neq p_2$.}{}
\qitem{Review · Factorise}{DEV}{Factorise $6x+9$.}{$3(2x+3)$.}{}
\qitem{Thread · Method}{DIV}{Write $360$ as a product of prime factors. How many divisors does it have?}{$360=2^3\times3^2\times5$; $(3+1)(2+1)(1+1)=24$ divisors.}{}
\end{fiche}
